\documentclass[11pt, a4paper]{article}

% load packages and settings
\input{packages.tex}

\title{Supplementary Appendix for: CD4-staged HIV back-calculation model including migration}
\author{}
\date{}

\begin{document}

\maketitle
\tableofcontents
\cleardoublepage

\section{Supplementary methods}

\subsection{Background of back-calculation models}

CD4-staged Bayesian `back-calculation' models, which account for HIV progression and time-varying probability of diagnosis, allow reconstruction of the underlying undiagnosed prevalence and incidence of infection from diagnosis data.

In the UK, data concerning HIV and AIDS diagnoses are collected centrally via epidemiological case surveillance, resulting in time-series of observed counts of quarterly HIV and AIDS diagnoses. Provided an estimate of the incubation period (i.e.\ the time from exposure to symptoms) is available, this can be combined with the observed data on diagnoses to `back-calculate' the incidence of infection. The general idea of back-calculation is that the time of symptom onset for an infected person is equal to the sum of the time of exposure and the incubation period~\parencite{egan_review_2015}. In continuous time $d(t)$, the rate at which new diagnoses occur at time $t$, can be expressed by the following convolution:
%
\[
  d(t) = \int_{t_0}^t{h(s)f(t-s) \dif s}
\]

where $d(t)$ is the rate of diagnosis at time $t$, $t_0$ is the starting time of the epidemic, $h(s)$ is the rate of new infections at time $s$, and $f(t-s)$ is the incubation period distribution function evaluated at time $t-s$~\parencite{brookmeyer_aids_1994}.

The occurrence of infections over time is typically modelled as a non-homogeneous Poisson process, with the number of new infections in a particular time interval being i.i.d. Poisson random variables with mean $h(s)$~\parencite{becker_method_1991, rosenberg_backcalculation_1991}. Extensions to the back-calculation model have incorporated HIV progression estimates based on CD4 counts at diagnosis~\parencite{sweeting_bayesian_2005, birrell_estimating_2012}, and age-specific incidence~\parencite{brizzi_extending_2019}.

To enable estimation of HIV incidence among groups with a large proportion of migrants we propose incorporating a migration process as part of the existing CD4 back-calculation model is necessary to enable , and will contribute to the estimation of pre- and post-migration incidence. This development requires detailed information on the country of acquisition for newly diagnosed individuals, with patterns of migration incorporated to inform recent trends.

% \subsection{Aims}

% \begin{enumerate}
%     \item To fit the CD4 including migration model to simulated data, where incidence of infection and migration are pre-defined quantities, propagated through the model dynamic equations, and can be compared to the estimates reconstructed by the model. 
%     \item To assess the additional benefit of considering migration in the CD4 back-calculation model, compared to a model which does not consider migration, and the feasibility of fitting the model to real data.
% \end{enumerate}

\section{Model description}\label{sec:model-description}

Several iterations of back-calculation models exist:
%
\begin{itemize}
    \item The CD4-staged back-calculation model~\parencite{birrell_estimating_2012}
    \item An age-dependent extension to the CD4-staged back-calculation model~\parencite{brizzi_extending_2019}
    \item A dual-biomarker (CD4 and RITA) back-calculation model (age-dependent and age-independent) (citation)
\end{itemize}

The model described here is an extension to the CD4-staged back-calculation model to include migration components. The addition of migration components will enable this model to be adapted to incorporate the developments listed above, i.e.\ age-dependency and inclusion of RITA information.

\subsection{Model dynamics}

Figure~\ref{fig:back-calc-migration} shows the model states and transition intensities for the CD4 back-calculation model including migration. UK incidence occurs subject to a time-varying intensity $(1 + p_j)h_j$, and migrants arrive in the country subject to a time-varying intensity $\omega_j$.
%
\begin{figure}[htbp!]
  \centering
  \small
  \begin{tikzpicture}
    % Nodes
    \node[] (m0) at (0,0) {};
    
    \node[rstate, align=center] (e1) [right = 20mm of m0] {CD4$\geq$500\\($e_{1,j}$)};
    \node[rstate, align=center] (e2) [right = 11mm of e1] {CD4 350-499\\($e_{2,j}$)};
    \node[rstate, align=center] (e3) [right = 11mm of e2] {CD4 200-349\\($e_{3,j}$)};
    \node[rstate, align=center] (e4) [right = 11mm of e3] {CD4<200\\($e_{4,j}$)};

    \node[state, align=center] (d1) [below = of e1] {CD4$\geq$500\\($\mu_{1,j}$)};
    \node[state, align=center] (d2) [below = of e2] {CD4 350-499\\($\mu_{2,j}$)};
    \node[state, align=center] (d3) [below = of e3] {CD4 200-349\\($\mu_{3,j}$)};
    \node[state, align=center] (d4) [below = of e4] {CD4<200\\($\mu_{4,j}$)};
    \node[state, align=center] (d5) [right = 11mm of e4] {AIDS\\($\mu_{5,j}$)};

    \node[] (i1) [above = of e1] {};
    \node[] (i2) [above = of e2] {};
    \node[] (i3) [above = of e3] {};
    \node[] (i4) [above = of e4] {};
    \node[state, align=center, color = Plum] (m1) [above right = 15mm and -4mm of e2] {Transfers\\of care from\\abroad\\$(\mu^\text{Tr}_j)$};
    \node[] (i5) [above = of m1] {};
    
    % Edges
    
    \path (m0) edge node[above] {$(1+p_j)h_j$} (e1);
    \path (e1) edge node[] {$q_1$} (e2);
    \path (e2) edge node[] {$q_2$} (e3);
    \path (e3) edge node[] {$q_3$} (e4);
    \path (e4) edge node[] {$q_4$} (d5);

    \path (e1) edge[] node['] {$d_{1,j}$} (d1);
    \path (e2) edge[] node['] {$d_{2,j}$} (d2);
    \path (e3) edge[] node['] {$d_{4,j}$} (d3);
    \path (e4) edge[] node['] {$d_{5,j}$} (d4);

    \path (i1) edge[dotted, color = Plum] node[] {$\omega_j m_{1,j}$} (e1);
    \path (i2) edge[dotted, color = Plum] node[] {$\omega_j m_{2,j}$} (e2);
    \path (i3) edge[dotted, color = Plum] node[] {$\omega_j m_{3,j}$} (e3);
    \path (i4) edge[dotted, color = Plum] node[] {$\omega_j m_{4,j}$} (e4);
    \path (i5) edge[dotted, color = Plum] node[] {$\omega_j (1 - \sum_i m_{i,j})$} (m1);

  \end{tikzpicture}
  \caption{CD4-staged back-calculation model including migration transitions. Rounded boxes indicate latent states, square boxes indicate diagnosed states. Solid lines indicate HIV progression and transition from latent to diagnosed state; dotted lines indicate migration.}\label{fig:back-calc-migration}
\end{figure}

Denote as $\bm{e}_j = (e_{1,j},\dots,e_{4,j})$ the expected number of individuals in $(t_{j-1},t_j]$ in each of the latent states, and as $\bm{\mu}_j = (\mu_{1,j},\dots,\mu_{5,j})$ the expected number of new diagnoses in $(t_{j-1},t_j]$ in each of the diagnosed states. The model dynamics may then be expressed as a set of recursive equations which include the matrices $\mathbf{Q}_j$ and $\mathbf{D}_j$:
%
\begin{align*}
  \bm{e}_j   & = \mathbf{Q}_j^\mathsf{T} \bm{e}_{j-1} + {[(1 + p_j)h_j, 0, \ldots,0]}^\mathsf{T} + \omega_j\bm{m}_j \\
  \bm{\mu}_j & = \mathbf{D}_j^\mathsf{T} \bm{e}_{j-1}
\end{align*}

where $p_j$ is the ratio of (individuals born abroad):(individuals born in the UK) who acquire their HIV in the UK, defined as:
%
\[
    p_j = \frac{\text{incidence in the UK at time $j$ among individuals born abroad}}{\text{incidence in the UK at time $j$ among individuals born in the UK}}
\]

The expected number of HIV and AIDS diagnoses in $(t_{j-1},t_j]$, are denoted, respectively, as:
%
\[
  \mu_j^\text{H} = \sum_{i=1}^4{\mu_{i,j}} \: \text{ and } \:
  \mu_j^\text{A} = \mu_{5,j}
\]

and the expected number of transfers of care from abroad in $(t_{j-1},t_j]$ are expressed as:
%
\[
  \mu^\text{Tr}_j = \omega_j (1-\sum_{i=1}^4 m_{i,j})
\]

Our proposed model is a simplification of a model with an explicit year of UK arrival component. Whilst including the year of arrival would enable estimation of the number of years spent in the UK prior to diagnosis, it adds significant model complexity (see~\ref{sec:year-of-arrival}), requires the imputation of any missing year of arrival information, and would provide limited added value compared against better suited methods (ref). 

\subsection{Modelling year of UK arrival}\label{sec:year-of-arrival}

% \begin{figure}[h!]
%     \centering
%     \includegraphics[width=\linewidth]{figures/cbirth_yarr.png}
%     \caption{Completeness of year of arrival information for individuals born abroad, by probable route of infection, UK, 2001-2022.}
%     \label{fig:cbirth-yarr}
% \end{figure}

To illustrate a model which does include year of arrival, let $k$ index the year of arrival for individuals born outside the UK, and $j$ index the year of HIV acquisition, where year 1 is the year that the epidemic starts. This model would then seek to reconstruct the $k \times j$ matrix:
%
\[
\begin{array}{cc|ccccc}
& & \multicolumn{5}{c}{\text{$k$ = year of UK arrival}} \\
& & 0 & 1 & 2 & \ldots & T_k \\
\hline
\multirow{3}{*}{\text{$j$ = year of incidence}} & 1 & h_1 & h_{1,1} & h_{1,2} & \ldots & h_{1,T_k} \\
& 2 & h_2 & h_{2,1} & h_{2,2} & \ldots & h_{2,T_k} \\
& \vdots & \vdots & \vdots & \vdots & \ddots & \vdots \\
& T_j & h_{T_j} & h_{T_j,1} & h_{T_j,2} & \ldots & h_{T_j,T_k} \\
\end{array}
\]

The first column, $h_j$, represents individuals born in the UK who acquire HIV in the UK at time $j$, having been either born in the UK, or arriving prior to the start of the HIV epidemic. The entries below the diagonal, $\{h_{j,k};\ j>k\}$, represent individuals who acquire their HIV after UK arrival, whereas the entries above (or on) the diagonal, $\{h_{j,k};\ j \leq k\}$, represent individuals who acquire their HIV prior to UK arrival. For simplicity, let $T_k \equiv T_j$, and assume any individuals who arrived before the start of the HIV epidemic acquired their HIV in the UK.

A key challenge with estimating this matrix is that each additional year of data, $j$, increases the matrix size by $2j$ and requires both vertical and horizontal summation to obtain parameters of interest:
%
\[
    p_jh_j = \sum_{k=1}^{j-1} h_{j,k} \quad \quad \quad
    \omega_k = \sum_{j=1}^{k} h_{j,k}
\]

The calculation of $\omega_k$ can be simplified by not explicitly modelling incidence outside of the UK, but this will still involve $j$ additional matrix components to be estimated for each additional year of data.

Instead, in the proposed model, we consider only two columns: UK-born, and born abroad. As before, individuals who arrived before the start of the HIV epidemic are assumed to have acquired their HIV in the UK, and included in the UK-born column:
%
\[
\begin{array}{cc|cc}
& & \text{UK-born} & \text{born abroad} \\
\hline
\multirow{3}{*}{\text{$j$ = year of incidence}} & 1 & h_1 & p_1h_1 + \omega_1 \\
& 2 & h_2 & p_2h_2 + \omega_2 \\
& \vdots & \vdots & \vdots \\
& T & h_T & p_Th_T + \omega_T \\
\end{array}
\]

The born abroad column includes all individuals who arrived in the UK after the start of the HIV epidemic, with the $j$th entry being the sum of: 
%
\begin{enumerate}
    \item[$p_jh_j$:] individuals born abroad with HIV acquired in the UK at time $j$, and,
    \item[$\omega_j$:] individuals born abroad with HIV acquired abroad who arrived at time $j$.
\end{enumerate}

\subsection{Statistical inference}

In this model, individuals are either present in the country at the point of HIV acquisition, arrive at a given time with undiagnosed HIV acquired abroad, are diagnosed with HIV or AIDS, or are transfers of care from abroad.  Hence individuals may be in one of four types of state in each time interval $(t_{j-1}, t_j]$:
%
\begin{enumerate}
  \item undiagnosed;
  \item diagnosed with HIV;
  \item diagnosed with AIDS;
  \item transfer of care from abroad.
\end{enumerate}

Assume HIV and AIDS diagnoses and transfers of care from abroad observed during quarter $j$ are independent Poisson random variables $Y^\text{H}_j$, $Y^\text{A}_j$, and $Y^\text{Tr}_j$ with expectations $\mu_j^\text{H}$, $\mu_j^\text{A}$, and $\mu_j^\text{Tr}$, respectively:
%
\begin{align*}
  Y^\text{H}_j & \sim \text{Poisson}(\mu_j^\text{H}) \\
  Y^\text{A}_j & \sim \text{Poisson}(\mu_j^\text{A}) \\
  Y^\text{Tr}_j & \sim \text{Poisson}(\mu_j^\text{Tr})
\end{align*}

Additionally, let $Y^\text{C}_{j}$ be multinomial random variables for the total number of diagnoses in each CD4 state, with expectations $\bm{\mu}_{j}^\text{C}$:
%
\begin{align*}
  Y^\text{C}_{j} & \sim \text {Multinomial}(\bm{\mu}_j^\text{C},n_j)
\end{align*}

and with parameters:
%
\begin{align*}
  \bm{\mu}^\text{C}_j =\left(\frac{\mu_{1,j}}{n_j},\frac{\mu_{2,j}}{n_j},\frac{\mu_{3,j}}{n_j},\frac{\mu_{4,j}}{n_j}\right) \quad \quad \quad \quad \quad
  n_j & = \sum_{i=1}^4{y_{i,j}}
\end{align*}

where $y_{i,j}$ are the observed data on HIV diagnoses in each of the four CD4 states, respectively, in $(t_{j-1},t_j]$.

\subsubsection{Likelihood}

Likelihood contributions are obtained from the observed data on HIV and AIDS diagnoses and transfers of care ($y_j^\text{H}$, $y_j^\text{A}$, and $y_j^\text{Tr}$) which include all diagnoses and transfers of care (including those with missing CD4 information) ($L_1$).
%
\begin{align*}
  L_1(\bm{y^\text{H}}, \bm{y^\text{A}}, \bm{y^\text{Tr}};\bm{h},\bm{d})                      & \propto \prod_{j=1}^T {(\mu_j^\text{A})}^{y^\text{A}_j}\exp(-\mu_j^\text{A}) \times {(\mu_j^\text{H})}^{y^\text{H}_j}\exp(-\mu_j^\text{H}) \times {(\mu_j^\text{Tr})}^{y^\text{Tr}_j}\exp(-\mu_j^\text{Tr})
\end{align*}

Likelihood contributions are also obtained from observed CD4 count data ($L_2$), conditional on $y_{i,j}$\@:
%
\begin{align*}
  L_2(\bm{y}^\text{C} \mid \bm{y}^\text{H};\bm{h},\bm{d}) & \propto \prod_{j=1}^T \prod_{i=1}^4 {(\mu_{i,j})}^{y_{i,j}} 
\end{align*}

Here bold lower-case letters represent vectors, i.e.\ $\bm{y}^\text{H} = (y_1^\text{H}, y_2^\text{H}, \ldots, y_T^\text{H})$. The full likelihood is the product of the two likelihood terms:
%
\begin{align*}
  L(\bm{y}^\text{H}, \bm{y}^\text{A},\bm{y}^\text{Tr};\bm{h},\bm{d}) = L_1(\bm{y^\text{H}}, \bm{y^\text{A}}, \bm{y^\text{Tr}};\bm{h},\bm{d}) \times L_2(\bm{y}^\text{C} \mid \bm{y}^\text{H};\bm{h},\bm{d})
\end{align*}

Finally, whilst not strictly a contribution to the likelihood, an approximate Markov melding approach~\cite{goudie_joining_2019} is used to inform $p_j = \exp(\eta_j)$:
%
\begin{align*}
    \log(\hat{p}_j) &\sim \text{Normal}(\eta_j, (\sigma^p)^2)
\end{align*}

Where $\hat{p}_j$ and $\sigma^{p}$ are the mean and standard deviation inferred from a previous study of country of birth among individuals who acquire their HIV in the UK, and $\eta_j$ is defined via a logit transformation:
%
\begin{align*}
    \eta_j &= \text{logit}\left(\frac{p_j}{1+p_j}\right) \\
        &= \log\left[\left(\frac{p_j}{1+p_j}\right) / \left(1 - \frac{p_j}{1+p_j}\right)\right]\\
        &= \log(p_j)
\end{align*}

\subsubsection{Smoothing specification using splines and random walks}

For identifiability, the time-varying infection ($\bm{h}$), migration ($\bm{\omega})$, diagnosis probabilities ($\bm{d}$), country of birth ratio ($\bm{p}$), and migration probabilities ($\bm{m}$) are specified according to splines on the logit-scale or log-scale. Priors on the hyperparameters for each of these splines were chosen to reflect prior beliefs, or left as uninformative. This specification serves to smooth these processes over time, ensuring that sudden jumps do not occur between time intervals~\parencite{birrell_estimating_2012, brizzi_estimating_2017}. Details of the specific spline and prior specification for each process are included below.

\textbf{Infection process}

The infection process, $h_j$, is modelled by:
%
\[
h_j = \exp(\lambda_j) \quad\quad
\lambda_j = \mathbf{X}^{\text{inf}}_j \boldsymbol{\beta}^{\text{inf}}
\]

where $\mathbf{X}^{\text{inf}}_j$ is the thin plate regression spline basis matrix with $k=10$ knots and $m=2$ penalty order, with priors:
%
\[
\beta^{\text{inf}}_1 \sim \mathcal{N}(5, 1.5^2) \quad\quad
\beta^{\text{inf}}_{2:10} \sim \mathcal{N}(0, (\sigma^{\text{inf}})^2) \quad\quad
\log \sigma^{\text{inf}} \sim \mathcal{N}(-2, 0.5)
\]

\textbf{Migration process}

The migration process, $\omega_j$ is modelled by:
%
\[
\omega_j = \exp(\gamma_j) \quad\quad
\gamma_j = \mathbf{X}^{\text{mig}}_j \boldsymbol{\beta}^{\text{mig}}
\]

where $\mathbf{X}^{\text{mig}}_j$ is the thin plate regression spline basis matrix with $k=10$ knots and $m=2$ penalty order, with priors:
%
\[
\beta^{\text{mig}}_1 \sim \mathcal{N}(4, 1) \quad\quad
\beta^{\text{mig}}_{2:10} \sim \mathcal{N}(0, (\sigma^{\text{mig}})^2) \quad\quad
\log \sigma^{\text{mig}} \sim \mathcal{N}(-2, 0.5)
\]

\textbf{Diagnosis probabilities}

The diagnosis probabilities, $d_{i,j}$ are modelled by:
%
\[
d_{i,j} = \text{logit}^{-1}(\delta_{i,j}) \quad\quad 
\delta_{i,j} = \alpha^{\text{diag}}_i + \mathbf{w}_i \cdot \mathbf{X}^{\text{diag}}_j \boldsymbol{\beta}^{\text{diag}}_i \quad\quad
i \in \{1, \ldots, 4\}
\]

where $\mathbf{X}^{\text{diag}}_j$ is the thin plate regression spline basis matrix with $k=6$ knots and $m=2$ penalty order, and $\mathbf{w}_i = [0.5, 0.5, 0.5, 0.7]$ is a weighting term, with priors:
%
\[
\alpha^{\text{diag}}_i \sim \mathcal{N}(-3, 0.5) \quad\quad
\beta^{\text{diag}}_{i,1} \sim \mathcal{N}(0, 1) \quad\quad
\beta^{\text{diag}}_{i,2:6} \sim \mathcal{N}(0, (\sigma^{\text{diag}}_i)^2) \quad\quad
\sigma^{\text{diag}}_i \sim \text{Gamma}(2, 64)
\]

\textbf{Country of birth ratio}

The country of birth ratio, $p_j$ is modelled by:
%
\[
p_j = \frac{\text{logit}^{-1}(\eta_j)}{1 - \text{logit}^{-1}(\eta_j)} = \exp(\eta_j)\quad\quad
\eta_j = \alpha^{\text{cob}} + w^{\text{cob}} \cdot \mathbf{X}^{\text{cob}}_j \boldsymbol{\beta}^{\text{cob}}
\]

where $\mathbf{X}^{\text{cob}}_j$ is the thin plate regression spline basis matrix with $k=6$ knots and $m=2$ penalty order, and $w^{\text{cob}} = 0.5$ is a weighting term, with priors:
%
\[
\alpha^{\text{cob}} \sim \mathcal{N}(0, 1) \quad\quad
\beta^{\text{cob}}_1 \sim \mathcal{N}(0, 1) \quad\quad
\beta^{\text{cob}}_{2:6} \sim \mathcal{N}(0, (\sigma^{\text{cob}})^2) \quad\quad
\sigma^{\text{cob}} \sim \text{Gamma}(2, 64)
\]

\textbf{Migration probabilities}

The migration probabilities, $m_{i,j}$, are modelled by:
%
\[
\theta_{i,j} = \alpha^{\text{mp}}_i + \mathbf{X}^{\text{mp}}_j \boldsymbol{\beta}^{\text{mp}}_i \quad \text{for } i \in \{1, 2, 3, 5\}
\]

where $\mathbf{X}^{\text{mp}}_j$ is the thin plate regression spline basis matrix with $k=8$ knots and $m=2$ penalty order, with a sequential logit transformation to ensure the probabilities at time $j$ sum to 1:
%
\begin{align*}
m_{5,j} &= \text{logit}^{-1}(\theta_{5,j}) \\
m_{1,j} &= \text{logit}^{-1}(\theta_{1,j}) \cdot (1 - m_{5,j}) \\
m_{2,j} &= \text{logit}^{-1}(\theta_{2,j}) \cdot (1 - m_{5,j} - m_{1,j}) \\
m_{3,j} &= \text{logit}^{-1}(\theta_{3,j}) \cdot (1 - m_{5,j} - m_{1,j} - m_{2,j}) \\
m_{4,j} &= 1 - m_{5,j} - m_{1,j} - m_{2,j} - m_{3,j}
\end{align*}

and with priors for $i \in \{1, 2, 3, 5\}$:
%
\[
\alpha^{\text{mp}}_i \sim \mathcal{N}(0, 1) \quad\quad
\beta^{\text{mp}}_{i, 1} \sim \mathcal{N}(0, 1) \quad\quad
\beta^{\text{mp}}_{i, 2:8} \sim \mathcal{N}(0, (\sigma^{\text{mp}}_i)^2) \quad\quad
\sigma^{\text{mp}}_i \sim \text{Gamma}(2, 64)
\]

\subsubsection{Posterior predictive}

Similar post-hoc methods to those presented by~\cite{birrell_estimating_2012} are used to derive quantities of interest involving the posterior predictive distribution of the latent model states.

In the Bayesian framework, the priors are updated by the likelihood to obtain the posterior. As the posterior in this case is intractable, samples are obtained via a suitable MCMC algorithm (code for implementing this model in the STAN software is provided at~\url{https://pkirwan.github.io/cd4backcalc}).

\section{Simulation scenarios}\label{sec:simulation-scenarios}

A number of simulations will be performed to evaluate the performance of the model according to different model parameters and inform model choice.

\subsection{Data generating procedure}

Data are generated by propagating pre-defined time-series of quantities ($h_j$, $\omega_j$) and ratios/proportions ($p_j$, $m_{i,j}$, $d_{i,j}$) through the following convolution equations:
%
\begin{align*}
  \bm{e}_j   & = \mathbf{Q}_j^\mathsf{T} \bm{e}_{j-1} + {[(1 + p_j)h_j, 0, \ldots,0]}^\mathsf{T} + \omega_j\bm{m}_j \\
  \bm{\mu}_j & = \mathbf{D}_j^\mathsf{T} \bm{e}_{j-1}
\end{align*}

with $\mathbf{Q}_{j}$ and $\mathbf{D}_j$ defined as appropriate transition matrices for HIV progression and HIV diagnosis.

Incidence of UK-acquired infection over time ($h_j$) and migration of individuals with HIV acquired abroad ($\omega_j$) will be drawn from the following distributions:
%
\begin{align*}
  h_j &\sim \text{Poisson}(\lambda_j) \\
  \omega_j &\sim \text{Poisson}(\gamma_j)
\end{align*}

State-specific migration and diagnosis probabilities will be drawn from the following distributions, with all probabilities constrained to be positive and log/logistic transformation as appropriate:
%
\begin{align*}
  d_{i,j} &\sim \text{Multivariate Normal}(\delta_{i,j}, 0.01) \\
  p_j &\sim \text{Normal}(\eta_j, 0.01) \\
  m_{i,j} &\sim \text{Multinomial}(\omega_j, \theta_{i,j})
\end{align*}

HIV progression probabilities ($q_i$) will be fixed and non-time-varying:
%
\[
  q = \left[0.0976, 0.115, 0.116, 0.148\right]
\]

Initial prevalence in each of the latent states will be defined as:
%
\begin{align*}
  \text{initial prevalence UK-born, UK-acquired} &= h_1 * \left[4, 2, 1, 0.25\right] \\
  \text{initial prevalence born-abroad, UK-acquired} &= 0.2 * h_1 * \left[4, 2, 1, 0.25\right] \\
  \text{initial prevalence born-abroad, acquired abroad} &= \left[0, 0, 0, 0\right]
\end{align*}

After being propagated through the convolution equations the derived number of HIV, AIDS, and transfers of care from abroad will be rounded to integer values, with a ranked rounding algorithm which preserves totals across the CD4 strata.

\subsection{List of scenarios}

500 datasets will be generated for each scenario, described below and shown in graphical form in Figure~\ref{fig:sim-scenarios}.
%
\begin{enumerate}
  \item Constant (\texttt{constant\_scenario}): 
  \begin{itemize}
      \item $\lambda_j = 400$;
      \item $\gamma_j = 100$; 
      \item $\eta_j = 0.2$; 
      \item $\delta_{j} = \left[0.05, 0.08, 0.11, 0.25\right]$;
      \item $\theta_{j} = \left[0.2, 0.2, 0.2, 0.2, 0.2\right]$
  \end{itemize}

  \item Increasing incidence (\texttt{h\_increasing}), as for \texttt{constant\_scenario}, except:
  \begin{itemize}
      \item $\lambda_j = \text{seq}(200, 600)$;
  \end{itemize}

  \item Increasing leading to decreasing incidence (\texttt{h\_varying}), as for \texttt{constant\_scenario}, except:
  \begin{itemize}
      \item $\lambda_j = \text{seq}(750, 1500, 1100, 800, 200)$;
  \end{itemize}

  \item Increasing migration (\texttt{o\_increasing}), as for \texttt{constant\_scenario}, except:
  \begin{itemize}
      \item $\gamma_j = \text{seq}(50, 1000)$;
  \end{itemize}

  \item Spike in migration (\texttt{o\_step\_change}), as for \texttt{constant\_scenario}, except:
  \begin{itemize}
      \item $\gamma_j = \text{seq}(50, 600, 200, 300, 400, 2000)$;
  \end{itemize}

  \item Increasing diagnosis probabilities (\texttt{d\_varying}), as for \texttt{constant\_scenario}, except:
  \begin{itemize}
      \item $\delta_{j} = \left[\text{seq}(0.02,0.03, 0.01), \text{seq}(0.05,0.04, 0.08), \text{seq}(0.10,0.07, 0.07), \text{seq}(0.25,0.35, 0.30)\right]$;
  \end{itemize}

  \item Increasing ratio born abroad (\texttt{p\_increasing}), as for \texttt{constant\_scenario}, except:
  \begin{itemize}
      \item $\eta_j = \text{seq}(0.1, 0.4)$;
  \end{itemize}

  \item Different migration probabilities for each CD4 strata (\texttt{m\_differing}), as for \texttt{constant\_scenario}, except:
  \begin{itemize}
      \item $\theta_{j} = \left[0.1, 0.2, 0.2, 0.1, 0.4\right]$
  \end{itemize}

  \item Time-varying migration probabilities for each CD4 strata (\texttt{m\_varying}), as for \texttt{constant\_scenario}, except:
  \begin{itemize}
      \item $\theta_{j} = \left[\text{seq}(0.1,0.3), \text{seq}(0.2,0.22), \text{seq}(0.25,0.15), \text{seq}(0.3,0.05), \text{seq}(0.2, 0.4)\right]$
  \end{itemize}

  \item \texttt{combo\_1}: combine \texttt{h\_varying}, \texttt{o\_step\_change}, and \texttt{d\_varying}, with all other parameters kept constant.

  \item \texttt{combo\_2}: combine \texttt{h\_varying}, \texttt{o\_step\_change}, \texttt{d\_varying}, and \texttt{p\_increasing}, with all other parameters kept constant.

  \item \texttt{combo\_3}: combine \texttt{h\_varying}, \texttt{o\_step\_change}, \texttt{d\_varying}, \texttt{p\_increasing}, and \texttt{m\_varying}.

\end{enumerate}

\begin{figure}[h!]
    \centering
    \includegraphics[width=\linewidth]{figures/parameter_visualization_table.pdf}
    \caption{Graphical representation of the parameter values used in each simulation scenario.}
    \label{fig:sim-scenarios}
\end{figure}

\section{Prior information}\label{sec:prior-information}

As above, prior information on $\eta_j$ is included via approximate Markov melding:
%
\begin{align*}
    \log(\hat{p}_j) &\sim \mathcal{N}(\eta_j, (\sigma^p)^2)
\end{align*}

This information is essential for the model parameters to be well identified. Three prior specifications will be explored:
%
\begin{itemize}
    \item Model with no specification for $\hat{p}_j$ and $\sigma^p$
    \item Model with correctly specified $\hat{p}_j$ and $\sigma^p$
    \item Model with incorrectly specified $\hat{p}_j$ and $\sigma^p$
\end{itemize}

\subsection{Predictive mean squared error}

The  bias and predictive mean squared error (PMSE) of the posterior mean HIV incidence will be evaluated for each specification listed above. The bias is the expected difference between the parameter value specified for the data generating process, $\lambda_j$, and the estimated posterior mean incidence, $\hat{h}_{i,j}$, at each time point, $j = 1, \ldots, N$, over the 500 simulation iterations, $i = 1, \ldots, 500$:
%
\begin{align*}
  \text{Bias}(\hat{h}_j, \lambda_j) &= \frac{1}{500} \sum_{i=1}^{500} \hat{h}_{i,j} - \lambda_j
\end{align*} 

The PMSE is the expectation of the squared differences over the 500 simulation iterations:
%
\[
  \text{PMSE}(\hat{h}) = \frac{1}{500} \frac{1}{N} \sum_{i=1}^{500} \sum_{j=1}^N{(\hat{h}_{i,j} - \lambda_{i,j})}^2
\]

where $N$ is the total number of simulated time points. A lower PMSE indicates that the parameter values are more accurately estimated. The bias and PMSE of the posterior mean HIV incidence are reported for each scenario, and the distribution of the differences and squared differences over the 500 simulations is shown.

\cleardoublepage

\section{Supplementary figures}

\subsection{Comparison of simulated data to real data}

Figures~\ref{fig:quarterly-dx} and~\ref{fig:sim-dx} show the real diagnosis data extracted from the UK HIV database and the quantities generated for one simulation under the \texttt{combo\_3} scenario, demonstrating the similarity between the simulated values and real data.

\begin{figure}[h!]
    \centering
    \includegraphics[width=\linewidth]{figures/quarterly_diagnoses.png}
    \caption{Quarterly number of HIV and AIDS diagnoses, transfers from abroad, and HIV diagnoses by CD4 count, England, 1996 to 2023.}
    \label{fig:quarterly-dx}
\end{figure}

\begin{figure}[h!]
    \centering
    \includegraphics[width=\linewidth]{figures/sim_diagnoses.png}
    \caption{Simulated quarterly number of HIV and AIDS diagnoses, transfers from abroad, and HIV diagnoses by CD4 count from one simulation, for 120 quarters.}
    \label{fig:sim-dx}
\end{figure}

\clearpage

\subsection{Estimation of unknown parameters using correctly specified model}

The model was fitted to each of the 500 datasets, with 4 parallel chains and 3000 MCMC iterations. Trace plots were examined for divergent chains. Figures~\ref{fig:incidence-sim} to~\ref{fig:diag-prob-sim} show the results of these model fits, for each of the 12 scenarios.

\begin{figure}[h!]
    \centering
    \includegraphics[width=\linewidth]{figures/incidence_sim.png}
    \caption{Estimated incidence of HIV for each simulation scenario, compared to the simulated parameter value.}
    \label{fig:incidence-sim}
\end{figure}

\begin{figure}[h!]
    \centering
    \includegraphics[width=\linewidth]{figures/migration_sim.png}
    \caption{Estimated migration for each simulation scenario, compared to the simulated parameter value.}
    \label{fig:migration-sim}
\end{figure}

\begin{figure}[h!]
    \centering
    \includegraphics[width=\linewidth]{figures/migration_prob_sim.png}
    \caption{Estimated probability of migration by CD4 stratum for each simulation scenario, compared to the simulated parameter value.}
    \label{fig:migration-prob_sim}
\end{figure}

\begin{figure}[h!]
    \centering
    \includegraphics[width=\linewidth]{figures/prop_uk_sim.png}
    \caption{Estimated ratio of UK-born:born abroad for each simulation scenario, compared to the simulated parameter value.}
    \label{fig:prop-uk-sim}
\end{figure}

\begin{figure}[h!]
    \centering
    \includegraphics[width=\linewidth]{figures/diag_prob_sim.png}
    \caption{Estimated diagnosis probabilities by CD4 stratum for each simulation scenario, compared to the simulated parameter value.}
    \label{fig:diag-prob-sim}
\end{figure}

\clearpage

\subsection{Value of correctly specified prior information}

\begin{figure}[h!]
    \centering
    \includegraphics[width=\linewidth]{figures/bias.png}
    \caption{Bias for HIV incidence, migration, and the ratio of UK-born:born abroad with correct, misspecified, and without prior information.}
    \label{fig:bias}
\end{figure}

\begin{figure}[h!]
    \centering
    \includegraphics[width=\linewidth]{figures/bias_mp.png}
    \caption{Bias for the migration probabilities into each CD4 stratum, with correct, misspecified, and without prior information.}
    \label{fig:bias-mp}
\end{figure}

\begin{figure}[h!]
    \centering
    \includegraphics[width=\linewidth]{figures/pmse.png}
    \caption{Predictive mean squared error for HIV incidence, migration, and the ratio of UK-born:born abroad with correct, misspecified, and without prior information.}
    \label{fig:pmse}
\end{figure}

\begin{figure}[h!]
    \centering
    \includegraphics[width=\linewidth]{figures/pmse_mp.png}
    \caption{Predictive mean squared error for the migration probabilities into each CD4 stratum, with correct, misspecified, and without prior information.}
    \label{fig:pmse-mp}
\end{figure}

\clearpage

\subsection{Value of migration specification}

\begin{figure}[h!]
    \centering
    \includegraphics[width=0.9\linewidth]{figures/incidence_ai_spl.png}
    \caption{Estimates from model without migration compared to simulated parameter values from scenario \texttt{combo\_3}.}
    \label{fig:no-migration}
\end{figure}

\begin{figure}[h!]
    \centering
    \includegraphics[width=0.9\linewidth]{figures/combo_3.png}
    \caption{Estimates from model with migration compared to simulated parameter values from scenario \texttt{combo\_3}.}
    \label{fig:combo-3}
\end{figure}

\clearpage

\section{Preliminary results fitting the model to real data}\label{sec:real-data}

\subsection{Availability of country of birth information}

\begin{figure}[h!]
    \centering
    \includegraphics[width=\linewidth]{figures/cbirth_pexp.png}
    \caption{Availability of country of birth information in UK HIV diagnosis data, by year of UK diagnosis.}
    \label{fig:cbirth-pexp}
\end{figure}

\subsection{Traceplots}

\begin{itemize}
    \item Prior specification for MSM: $\hat{p}_j = \text{seq}\left(0.3, 1.5\right) \quad \quad \sigma^p = 1$
    \item Prior specification for heterosexual men: $\hat{p}_j = \text{seq}\left(9, 2\right) \quad \quad \sigma^p = 1$
    \item Prior specification for heterosexual women: $\hat{p}_j = \text{seq}\left(9, 2\right) \quad \quad \sigma^p = 1$
\end{itemize}

\begin{figure}[h!]
    \centering
    \includegraphics[width=\linewidth]{figures/lambda_msm.png}
    \caption{Traceplots for $\lambda_j$ and $p_j$, for the model fitted to data for MSM.}
    \label{fig:lambda-msm}
\end{figure}

\begin{figure}[h!]
    \centering
    \includegraphics[width=\linewidth]{figures/gamma_msm.png}
    \caption{Traceplots for $\gamma_j$, $m_{5,j}$, and the log posterior density for the model fitted to data for MSM.}
    \label{fig:gamma-msm}
\end{figure}

\begin{figure}[h!]
    \centering
    \includegraphics[width=\linewidth]{figures/lambda_hetmen.png}
    \caption{Traceplots for $\lambda_j$ and $p_j$, for the model fitted to data for heterosexual men.}
    \label{fig:lambda-hetmen}
\end{figure}

\begin{figure}[h!]
    \centering
    \includegraphics[width=\linewidth]{figures/gamma_hetmen.png}
    \caption{Traceplots for $\gamma_j$, $m_{5,j}$, and the log posterior density for the model fitted to data for heterosexual men.}
    \label{fig:gamma-hetmen}
\end{figure}

\begin{figure}[h!]
    \centering
    \includegraphics[width=\linewidth]{figures/lambda_hetwomen.png}
    \caption{Traceplots for $\lambda_j$ and $p_j$, for the model fitted to data for heterosexual women.}
    \label{fig:lambda-hetwomen}
\end{figure}

\begin{figure}[h!]
    \centering
    \includegraphics[width=\linewidth]{figures/gamma_hetwomen.png}
    \caption{Traceplots for $\gamma_j$, $m_{5,j}$, and the log posterior density for the model fitted to data for heterosexual women.}
    \label{fig:gamma-hetwomen}
\end{figure}

\clearpage

\subsection{Preliminary results}

\begin{figure}[h!]
    \centering
    \includegraphics[width=\linewidth]{figures/incidence_migration_msm.png}
    \caption{Posterior median (blue) and 95\% CrI estimates (grey) of HIV incidence, migration, UK-born:born abroad ratio, diagnosis probabilities and migration probabilities for MSM, by chain, England, 1995-2023.}
    \label{fig:incidence-migration-msm}
\end{figure}

\begin{figure}[h!]
    \centering
    \includegraphics[width=\linewidth]{figures/diagnoses_transfers_msm.png}
    \caption{Posterior median (blue) and 95\% CrI estimates (grey) of HIV diagnoses and number of transfers from abroad compared to observed data points, for MSM, by chain, England, 1995-2023.}
    \label{fig:diagnoses-transfers-msm}
\end{figure}

\begin{figure}[h!]
    \centering
    \includegraphics[width=\linewidth]{figures/incidence_migration_hetmen.png}
    \caption{Posterior median (blue) and 95\% CrI estimates (grey) of HIV incidence, migration, UK-born:born abroad ratio, diagnosis probabilities and migration probabilities for heterosexual men, by chain, England, 1995-2023.}
    \label{fig:incidence-migration-hetmen}
\end{figure}

\begin{figure}[h!]
    \centering
    \includegraphics[width=\linewidth]{figures/diagnoses_transfers_hetmen.png}
    \caption{Posterior median (blue) and 95\% CrI estimates (grey) of HIV diagnoses and number of transfers from abroad compared to observed data points, for heterosexual men, by chain, England, 1995-2023.}
    \label{fig:diagnoses-transfers-hetmen}
\end{figure}

\begin{figure}[h!]
    \centering
    \includegraphics[width=\linewidth]{figures/incidence_migration_hetwomen.png}
    \caption{Posterior median (blue) and 95\% CrI estimates (grey) of HIV incidence, migration, UK-born:born abroad ratio, diagnosis probabilities and migration probabilities for heterosexual women, by chain, England, 1995-2023.}
    \label{fig:incidence-migration-hetwomen}
\end{figure}

\begin{figure}[h!]
    \centering
    \includegraphics[width=\linewidth]{figures/diagnoses_transfers_hetwomen.png}
    \caption{Posterior median (blue) and 95\% CrI estimates (grey) of HIV diagnoses and number of transfers from abroad compared to observed data points, for heterosexual women, by chain, England, 1995-2023.}
    \label{fig:diagnoses-transfers-hetwomen}
\end{figure}

\clearpage

\addcontentsline{toc}{section}{References}
\printbibliography

\end{document}
